\section{Introduction}

Quantum entanglement is a fundamental form of quantum correlation that
can be stronger than any classical correlation~\cite{EPR35, B64,CHSH69,JNVWY20}.
%
It plays a crucial role in a myriad of areas such as quantum computing, quantum information,
quantum complexity, quantum cryptography, condensed matter physics, etc~\cite{HHHH09,NC10,Watrous18}.
%
Hence, comprehending both the capabilities and constraints of quantum entanglement stands
as a crucial research endeavor.
%
However, entanglement can also pose challenges in numerous applications, such as quantum
key distribution and quantum proof systems~\cite{Renner08,KMY03,OW16,HHJWY17}.
%
This raises the natural question of designing quantum channels that convert quantum states
into unentangled states. 
%
For the purpose of applications, such channel, also called \emph{disentangler}, $\Phi: \cH \to \cK\otimes\cK$ can be defined to satisfy two
conditions: (i) for any $\ket\psi\in\cK$, there is preimage $\ket\phi$, such that $\Phi(\phi)=\psi\otimes\psi$;
and (ii) for any density operator $\phi\in \cH$, $\Phi(\phi)$ is close to \emph{separable}. 

The \emph{quantum de Finetti} type theorems~\cite{christandl2007one,KR05,Renner08} provide examples of
disentanglers.
%
A quantum de Finetti theorem quantifies the closeness of a \emph{permutation-invariant} $\ell$-partite
quantum state, to $k$-partite separable states when all but $k$ subsystems are traced out.
%
A standard quantum de Finetti theorem reads
\begin{theorem}[Quantum de Finetti~\cite{KR05}]\label{thm:std_quantum_de_finetti}
  For every $d,\ell\ge k \in \mathbb{N}^+$, the channel $\Lambda \colon (\mathbb{C}^{d})^{\otimes \ell} \to (\mathbb{C}^d)^{\otimes k}$
  defined as $\Lambda(\rho) = \Tr_{\ell-k}(1/\ell! \sum_{\pi \in \Sym_\ell} \pi \rho \pi^{\dag})$ satisfies
  \begin{align*}
    \norm{\Lambda(\rho) - \int \ket{\psi}\bra{\psi}^{\otimes k} d\mu}_1 \le \frac{2 k d^2}{\ell}.
  \end{align*}
\end{theorem}
Note that the error bound scales at least\footnote{If instead of making the state permutation invariant,
we project it onto the symmetric subspace, which is a perfectly valid and efficient operation in the
quantum setting, then the dependence on $d$ in~\cref{thm:std_quantum_de_finetti} improves from $d^2$ to $d$.}
as $d/\ell$, and in this version of the quantum de Finetti theorem, the parameters are known to be essentially
tight.
%
Consequently, if each subsystem is composed of $n$ qubits, then obtaining
a non-trivial error bound requires at least $\ell \ge d = 2^{n}$
subsystems, making this channel impractical for many applications.
%
This is conjectured to be essentially the best you can achieve. In particular,
it is conjectured that for any disentangler, the input dimension will be
exponential in the output dimension~\cite{ABDFS08} to achieve that the output is always $\epsilon$ close in trace distance to some separable states for any constant $\epsilon<1$. 

\paragraph{Dimension Independent Disentangler from Unentanglement}
While the original disentangler conjecture remains widely open, in this work,
we show that there is an explicit, efficient (BQP), and \emph{dimension independent}
quantum disentangler for $k$-partite (output) system starting from a bipartite
unentangled system. More precisely, we prove

\begin{restatable}[Disentangler from unentanglement]{theorem}
{TheoChannel}\label{thm:channel}
  Let $d,\ell \ge k \in \mathbb{N}^+$.
  There is an efficient channel $\Lambda \colon (\mathbb{C}^{d})^{\otimes \ell}\otimes (\mathbb{C}^{d})^{\otimes \ell} \to (\mathbb{C}^d)^{\otimes k}$
  such that for any density operators $\rho_1, \rho_2 \in \mathbb{C}^{d\ell}$ there is a distribution $\mu$ on pure states $\ket\psi \in \C^d$ satisfying
  \begin{align*}
    \norm{\Lambda(\rho_1 \otimes \rho_2) - \int \ket{\psi}\bra{\psi}^{\otimes k} d\mu}_1 \le \tilde O\left(\left(\frac{k^3}{\ell}\right)^{1/4}\right).
  \end{align*}
  Furthermore, product states of the form $\rho_1 = \rho_2 = \ket{\psi}\bra{\psi}^{\otimes \ell}$ are
  mapped to $\ket{\psi}\bra{\psi}^{\otimes k}$.
\end{restatable}

In contrast to the de Finetti disentangler, our disentangler from unentanglement features
error parameters that are independent of the input dimension entirely! 
%
Subsequently, we discuss applications of~\cref{thm:channel} in testing product states and
the gap amplification in quantum proof systems, culminating in a near-optimal gap amplification for the $\QMA^+(k)$ class: Any improvement on this gap amplification would imply $\QMA^\R(k)=\NEXP$.\footnote{
We don't want to distract the readers by the issue about quantum states over real or complex numbers.
In many cases, quantum computation over reals captures that over complex numbers. However, to the best of the authors' knowledge, this is unclear in the context of $\QMA(2)$. We have to use $\QMA^\R(k)$ to denote the proof systems where the proofs are guaranteed to have real amplitudes.}
%
We also anticipate that our tool will find further applications beyond those discussed in
this paper.

\subsection{Super Product Test}

The \emph{product test} was designed to test if a state $\ket{\phi}$ is close to $k$-partite
product state, \ie $\ket{\phi} \approx \ket{\phi_1} \otimes \cdots \otimes \ket{\phi_k}$, given
two copies of $\ket{\phi}$.
%
This test involves applying a sequence of swap tests to each of the $k$ subsystems of the two
copies $\ket\phi$.
%
Clearly, if $\ket\phi$ is indeed a $k$-partite product state, all the swap tests accept with certainty.
%
On the other hand, if $\ket\phi$ is entangled across the $k$ subsystems, some swap test will reject
with a probability that depends on the amount of entanglement.
%
It can be argued that the product test is optimal for ensuring perfect completeness, i.e., accepting
product states with certainty~\cite{HM13}. 

Despite its utility and elegance, the product test has two limitations.
%
Firstly, it only provides a guarantee concerning its input $\ket{\phi}\otimes \ket{\phi}$ which 
are destroyed after the test, yielding a single classical bit as output.
%
Very often in applications, one also needs some extra certified input states $\ket\phi$ to manipulate in subsequent computations after
the test.
%
Secondly, and  probably more irritatingly, the product test always accepts with some 
constant probability (say $\ge 1/2$) no matter how far $\ket{\phi}$ is from
being $k$-partite product, i.e., it has poor soundness.
%
These limitations can be resolved if you have more than 2 copies of $\ket\phi$~\cite{Kada_2008, SY23}.
%
For instance, given $\ell$ copies of $\ket\phi$, then one can adapt the product test to
sequentially apply projections on to symmetric subspace on the first, second, and subsequent
subsystems of all the copies of $\phi$.
%
Intuitively, this should give us a stronger test whose analysis was left as an open problem in~\cite{HM13}.
%
Recently, She and Yuen~\cite{SY23} analyzed this higher order version of product test achieving improved soundness.
%
We restate this higher order product test as relying on some $\ell$ unentangled equal copies of 
$\ket{\psi}$ to deduce a $k$-partite product structure of the input state.
%
One can require something even stronger on the input to achieve what we call \emph{super product test}.

\begin{restatable}{lemma}{LemSupProd}\label{lem:super_product}
  The super product test on input $\ket{\psi} \otimes (\ket{\phi_1}\ldots\ket{\phi_k})^{\otimes \ell}$ accepts with probability
  \[
  \frac{\ell}{(\ell+1)} \cdot \abs{\braket{\psi}{\phi_1}\ldots\ket{\phi_k}}^2 + \frac{1}{(\ell+1)}\mper
  \]
\end{restatable}

This super product test focuses on determining whether a target state $\ket\psi$ is a product state or not.
%
In addition to the target state, there are $\ell$ copies of an already $k$-partite product state that come to help.
%
This test is very natural and simple, except it seems to ask too much of its inputs: To compare, the high-order
product test requires some copies of a state whereas \cref{lem:super_product} requires some copies of an already
$k$-partite product state of the form $\ket{\phi_1}\otimes \cdots \otimes \ket{\phi_k}$.
%
We claim the super product test is not really asking for too much because our disentangler channel effectively
``amplifies'' the number of unentangled systems.
%
In particular, we can rely on just two unentangled proofs to enforce a state close to 
$(\ket{\phi_1}\otimes \cdots \otimes \ket{\phi_k})^{\otimes \ell}$ by~\cref{thm:channel}. 
For simplicity, consider $k$ unentangled pairs of untangled proofs where the $i$\textsuperscript{th} pair applied \cref{thm:channel} yields $\ket{\phi_i}^{\otimes \ell}$. 
%
Then run the super product test on a target state $\ket{\psi}$ and the $\ell$ copies
of already product states from our disentangler.
%
Furthermore, note that it is very cheap to instead enforce a state close to
$(\ket{\phi_1}\otimes \cdots \otimes \ket{\phi_k})^{\otimes 2\ell}$, allowing us
to reserve the extra $\ell$ copies of $\ket{\phi_1}\otimes \cdots \otimes \ket{\phi_k}$
as once the super product test passes, they can be used in any other computations as
a very good proxy of $\ket\psi$.
%
With this combination, we achieve arbitrarily good soundness without requiring 
more than $2k$ unentangled states\footnote{Naturally, the $2k$ unentangled states
need to get larger in dimension to achieve better soundness.} while obtaining a
guarantee about the output, rather than having just a single classical bit of output. 

The above discussion on how we apply our disentangler is general, which we summarize in~\cref{sec:prop_testing}.

\subsection{A Gap Amplification for $\QMA^+(2)$ up to Criticality}

Next, we turn to the unentangled quantum proofs, the so-called $\QMA(2)$ class~\cite{KMY03}
and its variants.
%
First, we provide some background on this subject.

The complexity of $\QMA(2)$ was shown to be closely related to a variety
of quantum and classical computational problems, \eg determining if a
mixed state is entangled given its classical description, as well as,
various forms of classical polynomial/tensor optimization
(see~\cite{HM13} for a more comprehensive list).
%
Despite considerable interest and effort (\eg~\cite{DPS04,ABDFS08,BT09,Beigi10,BCY11,GNN12,SW12,P12,BH13,BH15,HNW17}),
we still only know the trivial complexity bounds $\QMA \subseteq \QMA(2) \subseteq \NEXP$.

Even the fact that $\QMA(2)$ admits strong gap amplification is non-trivial and remained open for
about $10$ years before the seminar work of Harrow and Montanaro~\cite{HM13}.
%
With \cref{thm:channel}, it is easy to give a new proof of this fact, which we discuss
in~\cref{sec:prop_testing}. 

A variant of $\QMA(2)$, denoted $\QMA^+(2)$, with proofs of nonnegative amplitudes
was introduced by Jeronimo and Wu in~\cite{JW23}.
%
The goal of this variant was to capture many properties of $\QMA(2)$ while having more
structure in order to obtain a greater understanding.
%
Indeed, they showed that $\QMA^+(2)=\NEXP$ by designing a $\QMA^+(2)$ protocol for a $\NEXP$-complete problem with a constant gap.
%
On the other end of their result is the observation that $\QMA^+(2)\subseteq\QMA(2)$
provided that the completeness-soundness gap of $\QMA^+(2)$ is a sufficiently large constant.
%
This makes $\QMA^+(2)$ an intriguing class to study since either (i) showing that $\QMA^+(2)=\QMA(2)$, via possibly a gap amplification approach for $\QMA^+(2)$,
would characterize the complexity of $\QMA(2)$, or (ii) showing $\QMA^+(2) \ne \QMA(2)$ would
give a better upper bound $\QMA(2) \subsetneq \NEXP$.
%

By virtue of the unentanglement assumption of $\QMA^+(2)$ and the product test~\cite{HM13},
$\QMA^+(2)$ admits some non-trivial gap amplification.
%
For example, a gap of $1/\poly(n)$ can be amplified to a constant gap in which the completeness
becomes $1-\exp(-\poly(n))$ and the soundness becomes some constant strictly less than $1$.
%
Recently, Bassirian, Fefferman and Marwaha~\cite{bassirian2023quantum}, building on~\cite{JW23},
curiously showed that $\QMA^+(1)=\NEXP$ also with a constant gap.\footnote{It is not clear that
their gap can be made as large as the one for $\QMA^+(2)=\NEXP$.}
%
Since in the large constant gap regime of $\QMA^+(1)$, we have $\QMA^+(1)=\QMA \subseteq \textup{PP}$,
their result rules out the strong gap amplification for $\QMA^+(1)$ unless $\NEXP\subseteq\textup{PP}$.
%
Moreover, it also suggests that strategies aimed at amplifying the gap for $\QMA^+(2)$ must rely on
the unentanglement assumption.
%
This is precisely where the tools like the product test or our disentangler become essential.

With our disentangler, we make progresses towards understanding of $\QMA^+(2)$ versus $\QMA(2)$.
%
In particular, our progresses can be summarized as two aspects with two motivating questions.

\begin{center}
{\it  Motivating question 1. How crucial is the nonnegative amplitudes assumption to obtain $\QMA^+(2)=\NEXP$?}
\end{center}
Regarding our first motivating question, we show that the nonnegative amplitudes assumption can be almost
completely removed by considering unentangled quantum proofs of almost general \emph{real} amplitudes.
%
More precisely, we show that to capture $\NEXP$ it suffices to have unentangled proofs of the form
$\ket{\psi} = \sqrt{a} \ket{\psi_+} + \sqrt{1-a} \ket{\psi_{-}}$ where
$\ket{\psi_+}$ has nonnegative amplitudes, $\ket{\psi_{-}}$ only has
negative amplitudes and $\abs{a-(1-a)} \ge 1/\poly(n)$ with $a \in [0,1]$.
%
In words, we require the proofs to have slightly more $\ell_2$-probability mass ($1/\poly(n)$ extra mass) either on
nonnegative or negative amplitudes.
%
We refer to the quantity $\abs{a-(1-a)}$ as the $\ell_2$-sign bias of $\ket{\psi}$.
%
We call the associated complexity class $\textup{almost-QMA}^\R(k)$ (formally defined in~\cref{def:almost_qma_k}).
Our main complexity result can be stated as follows.

\begin{restatable}{theorem}{TheoAlmostGenAmp}\label{thm:almost_gen_amp}
  $\NEXP = \textup{almost-QMA}^\R(k)$ with unentangled proofs of $\ell_2$-sign bias of\footnote{The letter $n$ represents the
  input size and $b(n)$ is any polynomial time computable function bounded from below by a polynomial, \ie by $1/n^c$ for
  some constant $c > 0$.} $b(n) \ge \poly(1/n)$ and $k=\poly(1/b(n))$.
\end{restatable}

We obtain the above result by investigating the other motivating question: Since the power of $\QMA^+(k)$
ranges from $\NEXP$ to $\QMA(k)$ depending on the gap,


\begin{center}{\it Motivating question 2.
  How much can we amplify the gap of $\QMA^+(k)$?}
\end{center}

We make significant progress addressing this question.
%
Specifically, we show that a even more relaxed version of $\QMA^+(3)$, featuring
a single proof with nonnegative amplitudes and the other two with general amplitudes,
equals $\NEXP$, with completeness $1-\exp(-\poly(n))$ and soundness $1/2+1/\poly(n)$.
%
At the first glance, this looks like a ``just so so'' gap amplification.
%
It is noteworthy that achieving a slightly improved soundness of $1/2-1/\poly(n)$ would
imply $\QMA^\R(3)=\NEXP$.
%
In particular, if $\QMA^\R(3) \ne \NEXP$, then there is a sharp phase
transition in the complexity around the gap of a half.

\begin{restatable}{theorem}{TheoQMAThree}\label{thm:qma3}
  $\NEXP = \textup{QMA}^+(3)$ with completeness $c=1-\exp(-\poly(n))$
  and soundness $s=1/2+1/\poly(n)$. Furthermore, we can assume a
  particular case of $\textup{QMA}^+(3)$ in which two unentangled
  proofs have arbitrary amplitudes whereas only one unentangled proof has
  nonnegative amplitudes.
\end{restatable}

\begin{figure}[ht]
  \captionsetup{width=.8\linewidth}
  \centering
\begin{tikzpicture}[scale=1]  % Increase the scale for a longer axis
    % Draw the x-axis
    \draw[|->, thick] (0,0) -- (10,0);
    \node[left] at (0,0) {$0$};
    \node[right] at (10,0) {$1$};

    % Draw a tick at 1/2
    \draw[thick] (5,-0.1) -- (5,0.1) node[above] {$\frac{1}{2}$};

    % Draw an interval around 1/2
    \draw[|<->|, yshift=-0.3cm] (4.75,0) -- (5.25,0);
    \node[below] at (5, -0.3) {$\frac{1}{\text{poly}(n)}$};

    % Labels on top of the intervals
    \node[above] at (2.5, 0.1) {$\NEXP$};
    \node[above] at (7.5, 0.1) {$\QMA^{\R}(3)$};
\end{tikzpicture}
\caption{Gap and the complexity regime of the particular version of $\QMA^+(3)$ from~\cref{thm:qma3}.
         A gap below $1/2-1/\poly(n)$ corresponds to $\NEXP$, whereas a gap above $1/2+1/\poly(n)$
         corresponds to $\QMA^\R(3)$, illustrating a sharp phase transition.}
\end{figure}


\subsection{Organization}

We introduce notations and review basic concepts and facts in~\cref{sec:prelim}.
%
In~\cref{sec:papo}, we present an efficient multipartite disentangler (like) channel from bipartite unentanglement.
%
This construction relies on new de Finetti type properties concerning the interplay between entanglement and symmetry
which we explore in~\cref{sec:de_finetti}.
%
In~\cref{sec:prop_testing}, we delve into the utility of our disentangler where we elaborate a generic framework in the context of property testing. As one example, we present a new proof that $\QMA(2)$ admits strong gap amplification.
%
The final two sections are devoted to design new tests and derive the main complexity results in this paper.
In~\cref{sec:super_swap}, we present the super swap and super product test which leverage
unentanglement to achieve much improved soundness than the well-known swap and product tests.
%
Finally, we provide protocols for $\NEXP$ in~\cref{sec:protocol} leading to the main complexity results of this paper,
\cref{thm:almost_gen_amp} and~\cref{thm:qma3}.
